\documentclass[11pt]{article}
\usepackage{fontspec}
\usepackage[french]{babel}
\usepackage{amsmath}
\usepackage{amssymb}
\usepackage{amsthm}
\usepackage{mathrsfs}
\usepackage{tikz-cd}
\usepackage{enumitem}

\newcommand{\printedpage}[1]{\marginpar{\footnotesize p.~#1}}
\newcommand{\msp}[1]{\marginpar{\footnotesize [#1]}}
\newcommand{\Ho}{\mathrm{Ho}}
\newcommand{\lc}{\mathrm{lc}}
\newcommand{\uFib}{\underline{\mathrm{Fib}}}

\begin{document}
\setcounter{footnote}{131}

\printedpage{160}
% head: A. Grothendieck : Dérivateurs, Ch. XII, Édition des Universités de Montpellier II et Paris VII
% foot: Version 30 mars 2003

\msp{page 247}

\textbf{Corollaire.} {\itshape Description du foncteur
\[
(\xi_1,\xi_2)\longmapsto\xi_1\times_S\xi_2 : \Ho^{\lc}_W(S)\times\Ho^{\lc}_W(S)\longrightarrow\Ho^{\lc}_W(S).
\]
Ce foncteur est bel et bien un produit cartésien dans} $\Ho^{\lc}_W(S)$ (\footnote{mais c'est un cas particulier des cas de $\Ho_W$ (si $\mathcal{M}$ a des produits binaires), en remplaçant $\mathcal{M}$ par $\mathcal{M}/S$.}), {\itshape où les produits finis existent} (\footnote{Dire aussi que les foncteurs $f^* : \Ho^{\lc}_W(S)\longrightarrow\Ho^{\lc}_W(S')$ (pour $f : S'\longrightarrow S$) commutent aux produits finis.}).

Pour prouver ce dernier point, posant
\[
\mathcal{H}=\Ho^{\lc}_W(S),
\]
il faut montrer que les foncteurs
\[
\begin{tikzcd}
\mathcal{H} \arrow[r, shift left, "\delta"] & \mathcal{H}\times\mathcal{H} \arrow[l, shift left, "\pi"]
\end{tikzcd}
\]
sont adjoints, avec
\[
\begin{array}{rclcrcl}
\delta & = & \mathrm{diag}_{\mathcal{H}} & : & \xi & \longmapsto & (\xi,\xi)\\
\pi & = & \text{`produit'} & : & (\xi_1,\xi_2) & \longmapsto & \xi_1\times_S\xi_2.
\end{array}
\]
On va définir les homomorphismes d'adjonction
\[
\begin{array}{rcllrcl}
\mathrm{id}_{\mathcal{H}} & \xrightarrow{\alpha} & \pi\delta & \text{i.e.} & \xi & \xrightarrow{\alpha_\xi} & \xi\times_S\xi\\
\mathrm{id}_{\mathcal{H}\times\mathcal{H}} & \xleftarrow{\beta} & \delta\pi & \text{i.e.} & (\xi_1,\xi_2) & \xleftarrow{\beta_{\xi_1,\xi_2}} & (\xi_1\times_S\xi_2,\xi_1\times_S\xi_2).
\end{array}
\]
Pour ceci, on décrit $\mathcal{H}$ comme
\[
\mathcal{H}=(\uFib_W S)(W_S)^{-1}
\]
et on note que $\delta$, $\pi$ proviennent de foncteurs adjoints au niveau de $\mathcal{H}_0=\uFib_W S$
\[
\begin{tikzcd}
\mathcal{H}_0 \arrow[r, shift left, "\delta_0"] & \mathcal{H}_0\times\mathcal{H}_0 \arrow[l, shift left, "\pi_0"]
\end{tikzcd},
\]
lesquels donnent lieu à $\alpha_0$, $\beta_0$
\[
\begin{array}{rcl}
\mathrm{id}_{\mathcal{H}_0} & \xrightarrow{\alpha_0} & \pi_0\delta_0\\
\mathrm{id}_{\mathcal{H}_0\times\mathcal{H}_0} & \xleftarrow{\beta_0} & \delta_0\pi_0.
\end{array}
\]

\printedpage{161}
% head: A. Grothendieck : Dérivateurs, Ch. XII, Édition des Universités de Montpellier II et Paris VII
% foot: Version 30 mars 2003

\msp{page 248}

On définit $\alpha$, $\beta$ par passage aux catégories de fractions pour $\alpha_0$, $\beta_0$. Les 2 compatibilités, étant vérifiées pour $\alpha_0$, $\beta_0$, le sont aussi pour $\alpha$, $\beta$. Q.e.d.

Je dégage ici le

\textbf{Lemme.} {\itshape Soit \begin{tikzcd}
\mathcal{X} \arrow[r, shift left, "f"] & \mathcal{Y} \arrow[l, shift left, "g"]
\end{tikzcd} un couple de foncteurs adjoint, et soient $\Sigma$, $\Sigma'$ des localiseurs dans $\mathcal{X}$, $\mathcal{Y}$ tels que $f(\Sigma)\subseteq\Sigma'$, $g(\Sigma')\subseteq\Sigma$, de sorte qu'on a un couple de foncteurs
\[
\begin{tikzcd}
\mathcal{X}\Sigma^{-1} \arrow[r, shift left, "\bar{f}"] & \mathcal{Y}\Sigma'^{-1} \arrow[l, shift left, "\bar{g}"]
\end{tikzcd}.
\]
Alors $\bar{f}$, $\bar{g}$ sont adjoints, les flèches d'adjonction étant déduites de celles pour $f$, $g$.}

\textbf{Corollaire.} {\itshape Soit $(\mathcal{M},W)$ une catégorie de modèles, supposons que dans $\mathcal{M}$ les produits binaires (resp. les produits finis) existent, et que pour $u_1 : X_1\longrightarrow X'_1$ et $u_2 : X_2\longrightarrow X'_2$, si $u_1,u_2\in W$, alors $u_1\times u_2 : X_1\times X_2\longrightarrow X'_1\times X'_2$ est dans $W$. Alors dans $\mathcal{M}W^{-1}$ les produits binaires (resp. produits finis) existent, et le foncteur localisation $\mathcal{M}\longrightarrow\mathcal{M}W^{-1}$ y commute.}

On applique le lemme à
\[
\begin{tikzcd}
\mathcal{M} \arrow[r, shift left, "\delta"] & \mathcal{M}\times\mathcal{M} \arrow[l, shift left, "\pi"]
\end{tikzcd}
\qquad \text{et } W,\ W\times W,
\]
et pour l'objet final,
\[
\begin{tikzcd}
\mathcal{M} \arrow[r, shift left, "e"] & e \arrow[l, shift left, "\bullet"]
\end{tikzcd}
\qquad\qquad \text{et } W,\ W_0=\{\mathrm{id}_e\}.
\]

\msp{page 249}

Je vais maintenant construire le foncteur
\[
i_* : \Ho_W(\Psi)\longrightarrow\Ho_W(\tilde{\Psi}).
\]
Il me faut d'abord prouver le

\textbf{Lemme 1.} {\itshape Soit $\mathcal{F}$ la sous-catégorie pleine de $\mathcal{M}(\Psi)$, catégorie des diagrammes \begin{tikzcd}[row sep=small, column sep=small]
 & X_1 \arrow[dl, "p_1"'] \\
S & \\
 & X_2 \arrow[ul, "p_2"]
\end{tikzcd} dans $S$ \texttt{[plutôt dans }$\mathcal{M}$\texttt{, sur }$S$\texttt{]}, formée de ceux pour lesquels $p_1$, $p_2$ sont dans Fib$_W$. Soit $W_{\mathcal{F}}$ le localiseur induit sur $\mathcal{F}$ par $W(\Psi)$ ($W$-équivalences terme à terme). Alors
\[
\mathcal{F}W_{\mathcal{F}}^{-1}\xrightarrow{\sim}\mathcal{M}(\Psi)W(\Psi)^{-1}\ (\overset{\text{déf}}{=}\Ho_W(\Psi)).
\]}

\printedpage{162}
% head: A. Grothendieck : Dérivateurs, Ch. XII, Édition des Universités de Montpellier II et Paris VII
% foot: Version 30 mars 2003

Nous allons prouver ce lemme plus bas. En l'admettant, pour construire $i_*$, il suffit de construire
\[
\mathcal{F}\longrightarrow\Ho_W(\tilde{\Psi})
\]
qui applique $W_{\mathcal{F}}$ en des isomorphismes. Je prendrai, bien sûr,
\[
\begin{gathered}
\begin{tikzcd}[ampersand replacement=\&, row sep=small, column sep=small]
 \& X_1 \arrow[dl] \& \& \& \& X_1 \arrow[dl] \& \\
S \& \& \longmapsto \& \& S \& \& X_1\times_S X_2 \arrow[ul] \arrow[dl] \\
 \& X_2 \arrow[ul] \& \& \& \& X_2 \arrow[ul] \&
\end{tikzcd}\\
\mathcal{F}\quad\longrightarrow\quad\Ho_W(\Psi).
\end{gathered}
\] % [sic] \Ho_W(\Psi) -> \Ho_W(\tilde{\Psi}), the image is a diagram of type \tilde{\Psi} and the functor to be constructed is \mathcal{F}\longrightarrow\Ho_W(\tilde{\Psi})
Vérifions qu'il transforme flèches de $W_{\mathcal{F}}$ en isomorphismes.

\msp{page 250}

\[
\begin{tikzcd}[row sep=small, column sep=small]
 & X_1\times_S X_2 \arrow[dl] \arrow[dd] & & X'_1\times_{S'}X'_2 \arrow[ll, "f_3"'] \arrow[dl] \arrow[dd] \\
X_1 \arrow[dd] & & X'_1 \arrow[ll, crossing over, "f_1"'] & \\
 & X_2 \arrow[dl] & & X'_2 \arrow[ll, "f_2"'] \arrow[dl] \\
S & & S' \arrow[ll, "f_0"'] \arrow[from=uu, crossing over] &
\end{tikzcd}
\]

Il faut prouver que si $f_0,f_1,f_2\in W$, alors $f_3\in W$. Mais considérons l'image inverse du carré cartésien $Q$ par $f_0 : S'\longrightarrow S$, d'où un carré cartésien $Q_{S'}$ et une factorisation
\[
\begin{tikzcd}
Q & Q_{S'} \arrow[l, "\tilde{f}"'] & Q' \arrow[l, "g"'] \arrow[ll, bend left, "f"]
\end{tikzcd}
\]
de $Q'\longrightarrow Q$. Comme $f_0\in W$, et que $X_1$, $X_2$, $X_1\times_S X_2$ sont dans $\uFib_W S$, les flèches horizontales réalisant $Q_{S'}\longrightarrow Q$ sont toutes dans $W$. Ceci, plus le fait que $f_1,f_2\in W$, impliquent que dans le morphisme $g$, les flèches $g_1$, $g_2$ sont dans $W$. Il en résulte que la flèche
\[
g_3 : X'_1\times_{S'}X'_2\longrightarrow(X_1\times_S X_2)\times_S S'
\]
l'est aussi (théorie du produit dans $\uFib_W S'$ et $\Ho^{\lc}_W(S)$), donc $f_3=\tilde{f}_3g_3\in W$, q.e.d.

\msp{page 251}

Notons qu'on a un isomorphisme canonique
\[
\alpha : i^*i_*\xrightarrow{\sim}\mathrm{id}_{\Ho_W(\Phi)}. \tag*{$(*)$}
\]
J'ai envie de prouver que les foncteurs
\[
\begin{tikzcd}
\Ho_W(\Phi) \arrow[r, shift right, "i_*"'] & \Ho_W(\tilde{\Phi}) \arrow[l, shift right, "i^*"']
\end{tikzcd}
\]
\printedpage{163}
% head: A. Grothendieck : Dérivateurs, Ch. XII, Édition des Universités de Montpellier II et Paris VII
% foot: Version 30 mars 2003
sont adjoints, et que $(*)$ est une des flèches d'adjonction. Il faut définir l'autre flèche d'adjonction
\[
\beta : i_*i^*\longleftarrow\mathrm{id}_{\Ho_W(\tilde{\Phi})}. \tag*{$(**)$}
\]
Mais pour ceci, je dois utiliser le

\textbf{Lemme 2.} Soit $\mathcal{G}$ la sous-catégorie pleine de $\mathcal{M}(\tilde{\Phi})$, formée des \begin{tikzcd}
X_1 \arrow[d, "p_1"'] & X_3 \arrow[l, "q_2"'] \arrow[d, "q_1"] \\
S & X_2 \arrow[l, "p_2"]
\end{tikzcd} tels que $p_1,p_2\in\mathrm{Fib}_W$ et soit $W_{\mathcal{G}}$ le localiseur induit par $W(\tilde{\Phi})$ ($W$-équivalences terme à terme). Alors l'inclusion $\mathcal{G}\longrightarrow\mathcal{M}(\tilde{\Phi})$ induit une équivalence
\[
\mathcal{G}W_{\mathcal{G}}^{-1}\xrightarrow{\sim}\mathcal{M}(\tilde{\Phi})W(\tilde{\Phi})^{-1}\overset{\text{déf}}{=}\Ho_W(\tilde{\Phi}).
\]

Je vais pour le moment admettre ce

\msp{page 252}

lemme, tout comme le lemme 1 (p. 249). Je vais considérer $i^*$, $i_*$ comme des flèches entre les categories $\mathcal{F}W_{\mathcal{F}}^{-1}$, $\mathcal{G}W_{\mathcal{G}}^{-1}$. Elles sont induites par les flèches % [sic] categories -> catégories, missing accent
\[
\begin{tikzcd}
\mathcal{F} \arrow[r, shift right, "i_{0\,*}"'] & \mathcal{G} \arrow[l, shift right, "i_0^*"']
\end{tikzcd}
\]
entre $\mathcal{F}$ et $\mathcal{G}$ eux mêmes, % [sic] eux mêmes -> eux-mêmes, missing hyphen
\[
\begin{tikzcd}
 & X_1 \arrow[dd] & X_3 \arrow[l] \arrow[dd] & & X_1 \arrow[dd] & \\
i_0^* :\quad \mathcal{G}\quad\longrightarrow\quad\mathcal{F} : & & & \longmapsto & & \\
 & S & X_2 \arrow[l] & & S & X_2 \arrow[l]
\end{tikzcd}
\]
\[
\begin{tikzcd}
 & X_1 \arrow[dd] & & & X_1 \arrow[dd] & X_1\times_S X_2 \arrow[l] \arrow[dd] \\
i_{0\,*} :\quad \mathcal{F}\quad\longrightarrow\quad\mathcal{G} : & & & \longmapsto & & \\
 & S & X_2 \arrow[l] & & S & X_2. \arrow[l]
\end{tikzcd}
\]
Ces deux foncteurs sont adjoints l'un de l'autre, et compatibles aux localiseurs $W_{\mathcal{F}}$, $W_{\mathcal{G}}$. Donc les foncteurs déduits de $i_0^*$, $i_{0\,*}$ par passage aux catégories de fractions sont également adjoints (cf. lemme p. 248).

Ainsi, modulo la démonstration des lemmes 1 et 2, on a prouvé ceci :

\printedpage{164}
% head: A. Grothendieck : Dérivateurs, Ch. XII, Édition des Universités de Montpellier II et Paris VII
% foot: Version 30 mars 2003

\msp{page 253}

\textbf{Théorème.} {\itshape Soit $(\mathcal{M},W)$ une catégorie de modèles ($W\subseteq\mathit{Fl}\,\mathcal{M}$) satisfaisant l'axiome de factorisation gauche (p. 240), donnant donc lieu à une bonne théorie des fibres $W$-homotopiques. Soient
\[
\begin{tikzcd}[row sep=small, column sep=small]
 & 1 & & & & 1 \arrow[dr] & \\
\Psi=0 \arrow[ur] \arrow[dr] & & {} \arrow[r, "i"] & {} & \tilde{\Psi}=0 \arrow[ur] \arrow[dr] & & 3 \\
 & 2 & & & & 2 \arrow[ur] &
\end{tikzcd}
\]
(immersions ouvertes dans} $\mathrm{Cat}${\itshape ), et considérons le foncteur restriction
\[
i^* : \Ho_W(\tilde{\Psi})\longrightarrow\Ho_W(\Psi).
\]
Ce foncteur admet un adjoint à droite
\[
i_* : \Ho_W(\Psi)\longrightarrow\Ho_W(\tilde{\Psi}).
\]
Ce foncteur $i_*$ se `calcule' ainsi : sa valeur en \begin{tikzcd}
X_1 \arrow[d, "p_1"'] & \\
S & X_2 \arrow[l, "p_2"]
\end{tikzcd} est canoniquement isomorphe (dans} $\Ho_W(\tilde{\Psi})${\itshape ) à \begin{tikzcd}
X_1 \arrow[d] & X_1\times_S X_2 \arrow[l] \arrow[d] \\
S & X_2 \arrow[l]
\end{tikzcd}, si on suppose que $p_1$ ou $p_2$ est une $W$-fibration. Dans le cas général il se calcule en factorisant $p_1$ en
\[
X_1\xrightarrow{i_1}\bar{X}_1\xrightarrow{p'_1}S,\qquad i_1\in W,\ p'_1\in\mathit{Fib}_W,
\]
et en prenant \begin{tikzcd}
\bar{X}_1 \arrow[d] & \bar{X}_1\times_S X_2 \arrow[l] \arrow[d] \\
S & X_2 \arrow[l]
\end{tikzcd}, ou en

\msp{page 254}

procédant de façon symétrique, ou en factorisant simultanement $p_1$ et $p_2$, et prenant % [sic] simultanement -> simultanément, missing accent
\[
\begin{tikzcd}
\bar{X}_1 \arrow[d] & \bar{X}_1\times_S\bar{X}_2 \arrow[l] \arrow[d] \\
S & \bar{X}_2. \arrow[l]
\end{tikzcd}
\]}

\end{document}
